
\section{Introduction}

This paper originates from a project to design and implement
a programming language for the implementation
of logical algorithms, as they occur for example in theorem provers
or interactive verification systems.
In programming, one often has a situation where a 
{\bf struct} (or record) type has a tree-like type definition:
The type of the first field is fixed, but the 
types of later fields depends on the values of the 
fields that occur before it. 
The simplest such case is \verb|std::optional|, which consists
of a boolean that indicates whether the value is there.
Second simplest case is \verb|std::variant<T1, ..., Tn>|, which contains
an integer indicating which type \verb|Ti| is currently present.
A third case appears in parsing, when one uses attribute grammars.
Symbols in an attribute grammar consist of a tag (specifying
which symbol is present), optionally followed by an attribute, whose
type depends on the symbol tag. For example if the symbol tag is 
string, identifier, or scan error, 
the attribute will have type \verb|std::string|. 
If the symbol tag is double, the attribute must be \verb|double|. 
The fourth case arises from implementation of logic: A propositional
formula consists of a logical operator, followed by either a string, a
single formula, two formulas, or a sequence of formulas,
dependent on the logical operator in front. 
Current static type systems do not support such type dependencies. 
After checking the first field, access to the following fields has take
place through operations that can fail at run time. 
The access of these fields can be dynamically type checked, but
it is beyond current static type systems. 

We want to design and implement type checking
procedures that can check such dependent 
types for realistic programs. In this paper, we present a design
that we believe will work in practice. 
It uses a novel type of tree-automata,
that we call \emph{deterministic multiset automata}. Their
distinguishing feature is that they are able to process multisets, 
up to a limit. The limit is a result of the finite number of states.
We use DMA-s to model array ranges.
If one uses an index variable $ i $ to change the values in an
array $ A $ with a for-loop, we will model $ A[ 0 \ldots i \rangle $ 
and $ A[ i \ldots \rangle $ as states of a DMA.

In Section~\ref{Sect_data_terms}, we introduce the data 
that our procedure works on. These are trees over a small
set of primitive types. Proper trees are constructed by a tuple
operator, and by an array operator. 
Our type system does not aim to analyze arithmetical properties
of numeric types, as done in \cite{JhaVaz2021ref}. 
DMA's can handle properties of form `$i$ is interval $ [2,5] $'
but not '$i$ is even'. Also, they don't support arithmetic,
so a DMA would not be able to determine that $ j := i + 1 $ 
implies that `$j$ will be in the interval $ [3,6]. $' 

During type checking, every function body is checked independently.
Type checking is always combined with overload resolution. 
The source code contains imprecise names, which can be mapped
to zero, one or any number of candidate implementations.
The typechecker checks which available implementations are 
guaranteed to succeed (whose preconditions are guaranteed to be met).
It removes the candidates whose preconditions are not met. 
If no implementation remains, the type checker reports an error. 
If more than one available implementation remains, the type checker
has to pick the most specific implementation. If no implementation
is most specific, the result is a type checking error (ambiguous overload).

Standard compilers perform type checking and intermediate code generation
at the same time.
Expressions are type checked, the overloads are resolved, and the 
expression is translated into intermediate code. 
(See \cite{WiSe2023})
For our type system, this is not possible, because
subtypes can change (disappear)
during loop iteration. As a consequence, we need to
use an intermediate
representation in which overloads are not resolved. Instead, 
for every imprecise statement, we
insert all overload candidates into the intermediate representation. 
Consider for example:
\begin{verbatim}
   x = 0
   while ...
      print(x)
         // There are two overloads printE (for even numbers),
         // and print (for arbitrary numbers). 
      x = ( ?succ, ( ?succ, x )).
\end{verbatim}
In order to resolve \verb| print(x) |into \verb| printE(x) |, 
one must know that all
operations on \verb|x| in the loop preserve evenness.
Some of the operations appear after the print statement in the code.
So, we will insert both versions of \verb|print| into the 
intermediate representation,
approximate all values that $ x $ can have using DMAs, and verify
that $ x $ is always even. In the end, both versions of
\verb|print| are eligible, but \verb|printE| is more specific,
so that it can be picked. 

Unfortunately, approximating values of variables is more
complicated than we sketched above. It is not sufficient
to collect all possible values of all variables at each program point
and map them into DMA states.
We need to collect combinations of possible values of variables. 
This is due to the fact that there may be dependencies between variables,
that would otherwise get lost in the process.
For example, if one compares two propositional formulas, one
typically starts by comparing the operators. If the operators are distinct,
the formulas are not equal. After that, one can do a case analysis
on the operator. If one formula has a binary operator on top,
one can safely access the left and right subformula in both formulas,
because one knows that the operators are equal. 
This information would be lost, if one would keep at each
program point only possible
values, instead of tuples of combinations of values.

In the next section, we define \emph{data terms}. 
Their definition is straightforward. They are built-up
from primitive types using a tuple constructor, and an array constructor. 

In the next section~\ref{Sect_DMA}, we define
deterministic multiset automata. 
They are a straightforward generalization of standard,
deterministic tree automata (DFTA in \cite{TATA2007})
to include multiset nodes. Evaluation is always bottom-up. 

In Section~\ref{Sect_intermediate}, we present an intermediate
representation. It uses flattened expressions, which
is quite common in compilers, see e.g.
\cite{LLVM2004}, \cite{Muchnick_CDI_1997}, or \cite{WiSe2023}.
Different from the standard approach, it uses a 
$ \statresolve $ statement, which for a given imprecise
statement, contains all candidate overloads, together with 
their preconditions, encoded as DMAs. 
The representation does not use flow graphs,
because the resolve statement is not naturally represented in
flow graphs. Our recursively defined format is slightly more
complicated than standard flow graphs, but much more flexible.
In Section~\ref{Sect_range_states}, we approximate states by 
mappings from variables to states of multiset automata. 
In addition, for every variable that can serve as an index of an array, 
a state keeps track of the values before the index and after the index
as separate multisets. The approximations are called 
\emph{range states} because they keep track of ranges in
arrays. We explain the effect of different statements
on range states. 

